% r03_h6w_a: proof of Theorem rot:thm:H6W, part 1: conventions and the bridge, the a priori bounds, King's flow on the axis lattice, the two one-lattice stability inputs, and the reduction to the small-field ratios.
\providecommand{\Ztor}{\ensuremath{\Zun^{\rm tor}}}
\providecommand{\hinf}{\ensuremath{h_\infty}}

The proof of \Cref{rot:thm:H6W} occupies the rest of this paragraph, up to the subsubsection on the forward bound. It is a one-lattice statement: only the axis lattice $\Xk=\Lk$ on the torus $T\defeq T_1=\DJ$ occurs, with King's standard block-spin flow at block factor $5$ and King's own corner-rooted blocks (the centred blocks of \Cref{rot:conv:C} play no role here), and neither the rotated lattice, nor the hybrid flow, nor \Cref{rot:thm:sandwich} is used. Throughout, $\fc=(J,\psi,n)$ is in the class of \Cref{rot:thm:CT}, $h_0\defeq-n\psi$, $G$ is a real smooth $P_J$-periodic $2$-form sampled in the frame of $\Xk$, $t\in\R$, and $r=r(k)$, $\delta_r=5^r\eps_k=5^{-\lfloor\vr k\rfloor}$ are the depth and the coarse spacing of \Cref{rot:lem:r}. The argument compares two ratios on the same lattice: the walled curvature ratio, which by the bridge of \Cref{rot:prop:CT} is King's torus-convention ratio at a wall that moves with $k$ by a constant of order $\eps_k$, and the torus-convention ratio at a frozen wall; the frozen family is Cauchy by King's renormalisation theorem read modulo the field-independent constant of its scalar-source sector, and the moving family stays close to the frozen one because a constant shift of the scalar source changes the effective action by a bounded amount per site, which the rate $|u_k|\le C\eps_k$ of \Cref{rot:prop:CT}(b) beats.

\subsubsection*{Conventions, the bridge and the a priori bounds}

\begin{lemma}[the moving and the frozen wall]\label{rot:lem:H6conv}
Write $\Ztor_k(t;h)\defeq\Ztor_{\Xk,T}(tG,h)$ for the torus-convention functional of \Cref{rot:prop:CT}(c) and
\[
h_k\defeq h_0-\tfrac12\Delta_T(\eps_k)\mathbf 1,\qquad\hinf\defeq h_0-\tfrac12\Delta_T(0)\mathbf 1,\qquad u_k\defeq h_k-\hinf=\tfrac12\big(\Delta_T(0)-\Delta_T(\eps_k)\big),
\]
a constant with $|u_k|\le C\eps_k$. Then for every $k$ and every real $t$
\[
r_k(t)\defeq\langle e^{tF(G)}\rangle_{\fc,\Xk}=\frac{\Zun_{\Xk,T}(tG,h_0)}{\Zun_{\Xk,T}(0,h_0)}=\frac{\Ztor_k(t;h_k)}{\Ztor_k(0;h_k)},\qquad s_k(t)\defeq\frac{\Ztor_k(t;\hinf)}{\Ztor_k(0;\hinf)} .
\]
The family $\{h_k\}_{k}\cup\{\hinf\}$ is bounded in every $C^N(T)$ norm, $\|h\|_{C^N}\le\|h_0\|_{C^N}+\Delta_*/2$; $\Xk/P_J$ is King's lattice on his torus \cite[(2.21)]{K1}, so the right side of the identity for $r_k$ is literally the object of \cite[Thms.~3.1, 3.4]{K1} with the scalar source $h_k$, and King's constants, which depend on the sources only through bounds on finitely many derivatives (\cite[pp.~664--665, 677]{K1}: ``the sources $g,h$ are bounded by $C$'', ``their derivatives are uniformly bounded''), are uniform over the family. The physical mass $m^2$ is the same in both conventions and is never changed.
\end{lemma}
\begin{proof}
The first identity is the walled-ratio display of \Cref{rot:def:Zun} at $g=0$, with $-w=h_0$ (see also \Cref{rot:lem:walledalg}(c)): the wall is the scalar source $h_0$ up to a constant that cancels in the ratio. The second is \Cref{rot:prop:CT}(c) at $\Lambda=\Xk$, whose prefactor does not depend on $G$ and cancels in the ratio. The existence of $\Delta_T(0)=\lim_{s\to0}\Delta_T(s)$ and $|u_k|\le C\eps_k$ are the rate clause of \Cref{rot:prop:CT}(b), and $|\Delta_T(\eps_k)|,|\Delta_T(0)|\le\Delta_*$ give the $C^N$ bound, the shifts being constants. No sign condition on the scalar source enters King's bounds: $\lambda>0$ makes every single-site measure finite for every real source.
\end{proof}

\begin{lemma}[a priori bounds on the ratios]\label{rot:lem:H6A}
Let $C_G\defeq\sup_kq_{\Xk}(G)/2<\infty$ (\Cref{rot:lem:q}). For every $k$, every real smooth $P_J$-periodic $h$ and every real $t$,
\[
1\le\frac{\Ztor_k(t;h)}{\Ztor_k(0;h)}\le e^{t^2q_{\Xk}(G)/2}\le e^{C_Gt^2},\qquad\Big|\frac{\Ztor_k(z;h)}{\Ztor_k(0;h)}\Big|\le\frac{\Ztor_k(\Re z;h)}{\Ztor_k(0;h)}\le e^{C_G(\Re z)^2}\quad(z\in\C),
\]
and the same in the free-boundary convention. Consequently, if a sequence of such ratios converges for every real $t$, it converges locally uniformly on $\C$ to an entire function, and all moments of $F(G)$ converge to the Taylor coefficients of the limit.
\end{lemma}
\begin{proof}
The torus-convention walled state is a BFS measure of King type in the sense of the paragraph after \Cref{rot:def:Zun}: the torus counterterm and Wick constant only change the Gaussian factor of the single-site measures and a constant. By \Cref{lem:s02:GD} with $F_{\Xk}(G)=\ip A{\ell_{\Xk}}$ (\Cref{rot:lem:CM}), $\langle e^{tF(G)}\rangle\le e^{t^2q_{\Xk}(G)/2}$; by Jensen's inequality and $\langle F(G)\rangle=0$ under the charge conjugation of \Cref{rot:def:tori}, $\langle e^{tF(G)}\rangle\ge1$. For complex $z$, $|e^{zF}|=e^{(\Re z)F}$. The functions $z\mapsto\Ztor_k(z;h)/\Ztor_k(0;h)$ are entire (the integral converges absolutely for every complex $z$ by the same Gaussian bound) and are bounded on every disc $|z|\le\mathsf R$ by $e^{C_G\mathsf R^2}$, uniformly in $k$; convergence on $\R$ then gives locally uniform convergence on $\C$ by Vitali's theorem \cite{Hor90}, and Cauchy's estimates give the moments.
\end{proof}

\subsubsection*{King's flow on the axis lattice}

For a lattice $X_{k'}$ and a depth $r'$ with $5^{r'}\eps_{k'}<1$, King's standard flow \cite[\S3]{K1} on $X_{k'}/P_J$ after $r'$ steps produces the unit lattice $T^{(r')}$ of physical spacing $\delta\defeq5^{r'}\eps_{k'}$, the small-field indicator $\chi_{r'}$ \cite[(3.2)]{K1}, on whose support $|\phi_{r'}|\le p(\delta)(\lambda\delta)^{-1/4}$, and the effective action $S^{(r')}_{k'}(tG,h)$ of \Cref{rot:def:eff}, here in the torus-convention reading for which \cite{K1} is written: its Gaussian parts (quadratic forms, $\ln N_{r'}$, $\ln Z_{r'}(A^{(r')})$) are defined at the mass $m^2$ without sources \cite[(2.13)--(2.16)]{K1}, and the sources enter only through the curvature vertex \cite[(3.20)]{K1} and the $h$-vertex \cite[(3.21)]{K1} of the perturbative part $E'_{r'}$, expanded to King's order $\bar n=13$ with a remainder. Put
\[
I^{(r')}_{k'}(t;h)\defeq\int\chi_{r'}\,e^{-S^{(r')}_{k'}(tG,h)}\,dA_{r'}\,d\phi_{r'},\qquad
\rho^{(r')}_{k'}(t;h)\defeq\frac{I^{(r')}_{k'}(t;h)}{I^{(r')}_{k'}(0;h)},
\]
and $I_k\defeq I^{(r(k))}_k$, $\rho_k\defeq\rho^{(r(k))}_k$ at the lattice's own depth. Two facts are used throughout: the unit lattice has $|T^{(r')}|=|T|\delta^{-3}$ sites (\cite[(2.15)]{K1} in unit-lattice counting); and $5^r\eps_k=5^{r+1}\eps_{k+1}=\delta_r$, so that $\Xk$ at depth $r$ and $X_{k+1}$ at depth $r+1$ live on the same unit lattice $T^{(r)}$ with the same $\chi_r$ and the same field bounds (\cite[p.~656]{K1}: ``since $L^k\eps=L^{k+n}\eps'$, the fields $A_k,\phi_k$ have the same bounds in each integral'').

\subsubsection*{The two stability inputs at the walls}

\begin{proposition}[King's upper bound on the axis lattice, at a wall]\label{rot:prop:H6S1}
For $h$ in the family of \Cref{rot:lem:H6conv}, every $k$ and every admissible depth $r'$, with $\delta=5^{r'}\eps_k$,
\[
e^{-C\delta^\sigma|T|}\,I^{(r')}_k(t;h)\le\Ztor_k(t;h)\le e^{C\delta^\sigma|T|}\,I^{(r')}_k(t;h)+B_k,\qquad B_k\defeq e^{-p(\delta)^2+C|T|}\eps_k^{-q},
\]
with $\sigma>0$, $C=C(t,G,h_0,|T|)$ locally uniform in $t$, and $q=q(t,G,h_0)<\infty$, all independent of $k$, of $r'$ and of $h$ in the family.
\end{proposition}
\begin{proof}
This is \cite[Thm.~3.1, (3.3)--(3.4)]{K1} for King's torus with the sources $(tG,h)$, with the additive large-field term in the form in which King uses it (\cite[(3.11)]{K1}, \cite[(4.7)]{K2}), at block factor $5$ on the torus $\DJ$ aligned with King's blocks: \Cref{thm:a06:BALA}(iii) with \Cref{hyp:k01:BALint}, in the torus convention for which it is printed, so that no bridge is crossed. Two readings are made, the same two as in \Cref{rot:prop:H3}: the uniformity over the bounded source family (\Cref{rot:lem:H6conv}); and the envelope. The factor $e^{C|T|}$ of the large-field term is the ultraviolet stability bound \cite[(2.23)]{K1}, which fails as printed when $\int h\ne0$ (\Cref{rot:rem:log}); what the inductive proof supplies at a wall is a bounded $h$-linear $c$-number per step, $e^{C|T|+C'(1+k)|\int h|}=e^{C|T|}\eps_k^{-q}$. The loss is harmless against $e^{-p(\delta_r)^2}$, since $p(\delta_r)^2$ grows at least quadratically in $k$ (\Cref{rot:lem:r}) while $\ln\eps_k^{-q}=qk\ln5$ is linear. The sources enter the proof through the vertices (3.20)--(3.21) only, with constants depending on $\sup|G|$, $\sup|\nabla G|$, $\sup|h|$, $\sup|\nabla h|$ (\cite[pp.~664--665, 677]{K1}).
\end{proof}

\begin{proposition}[the polynomial lower bound at the wall, torus convention]\label{rot:prop:H6S2}
For $h$ in the family of \Cref{rot:lem:H6conv} and every $k$: $\Ztor_k(0;h)\ge e^{-C_1|T|}\eps_k^{q_1}$, with $C_1$, $q_1$ independent of $k$ and of $h$ in the family.
\end{proposition}
\begin{proof}
This is \Cref{rot:thm:FS} on $\Lambda=\Xk$, $T=T_1$, for the bounded family ($\|h\|_\infty\le H\defeq n\|\psi\|_\infty+\Delta_*/2$; the sign of $h$ is not used there), read in the torus convention: the proof of \Cref{rot:thm:FS} moves the wall into the bare scalar coefficient, compares with a higher King point by Jensen's inequality and the monotonicity of $\langle\sum_x\eps^3|\phi|^2\rangle$ in that coefficient, and uses the source-free stability bounds \cite[(1.14)]{Ba1}, \cite[(1.3)]{Ba2} at the higher points together with the integrated counterterm identity and the perturbative bounds on the one-point Wick function. Each of these steps is available for the torus-convention model with the same constants: \cite[(1.14)]{Ba1} and \cite[(1.3)]{Ba2} are printed for it (\Cref{thm:a06:BALA}(iii)); its bare coefficient, Wick constant and vacuum constant differ from the free-boundary ones by $\Delta_T(s)$, $\Delta c_T(s)$, $\Delta E_T(s)$, which are bounded uniformly in $s$ (\Cref{rot:prop:CT}(b)) and shift the King points by bounded amounts only; and the perturbative coefficients of the torus one-point function differ from the infinite-lattice ones by image sums, exponentially small in $\ell(T)$. The bound reads $\ln\Ztor_k(0;h)\ge-C|T|-C(1+k)|T|-qk\ln5$, which is the claim with $q_1=q+C|T|/\ln5$.

At the mass-increasing wall the same conclusion also follows from Jensen's inequality as in \Cref{rot:prop:J}, in the torus convention: $\Ztor_k(0;h)\ge\Ztor_k(0;0)\exp[\sS_{\Xk}(h)\,w^{\rm tor}_k]$ with $w^{\rm tor}_k\defeq\langle|\phi(x)|^2\rangle^{\rm tor}_{0,0}-2c_T(\eps_k)$, independent of $x$ by translation invariance ($P_J\subset\Xk$), where $\Ztor_k(0;0)\ge e^{-C|T|}$ is \cite[(1.14)]{Ba1} and $\sS_{\Xk}(h)=-n\sS_{\Xk}(\psi)-\frac12\Delta_T|T|<0$ by the wall condition of \Cref{rot:thm:CT}; together with an upper bound $w^{\rm tor}_k\le C(1+k)$ on the source-free one-point Wick function, which is not established here and not needed, this gives the claim with $q_1=C|\sS_{\Xk}(h)|/\ln5$. This is the only place in the proof where the sign of the wall would enter; the route through \Cref{rot:thm:FS} does not use it.
\end{proof}

\subsubsection*{The ratios are small-field ratios up to $o(1)$}

\begin{lemma}[the large-field term is negligible]\label{rot:lem:H6stepA}
Let $h$ be in the family of \Cref{rot:lem:H6conv}, $r'$ an admissible depth for $\Xk$, $\delta=5^{r'}\eps_k$, $x\defeq C\delta^\sigma|T|$ with the constant of \Cref{rot:prop:H6S1}, and $b_k(h)\defeq2B_k/\Ztor_k(0;h)$. Then
\[
b_k(h)\le2e^{-p(\delta)^2+(C+C_1)|T|}\eps_k^{-q-q_1},
\]
and whenever $b_k\le\frac12$ and $2x\le\ln2$:
\begin{equation}\label{rot:eq:H6stepA}
e^{-2x}(1-b_k)\,\rho^{(r')}_k(t;h)\le\frac{\Ztor_k(t;h)}{\Ztor_k(0;h)}\le e^{2x}\rho^{(r')}_k(t;h)+b_k,\qquad
\rho^{(r')}_k(t;h)\le4e^{C_Gt^2},\qquad
\Big|\frac{\Ztor_k(t;h)}{\Ztor_k(0;h)}-\rho^{(r')}_k(t;h)\Big|\le4e^{C_Gt^2}(e^{2x}-1)+(1+4e^{C_Gt^2})\,b_k .
\end{equation}
At the lattice's own depth $r=r(k)$, $x=C\delta_r^\sigma|T|\to0$ and $b_k(h)\to0$ as $k\to\infty$ (\Cref{rot:lem:r}: $p(\delta_r)^2$ grows at least quadratically in $k$, the exponent of $\eps_k^{-q-q_1}$ linearly), so $|r_k(t)-\rho_k(t;h_k)|\to0$ and $|s_k(t)-\rho_k(t;\hinf)|\to0$ for every real $t$, locally uniformly in $t$.
\end{lemma}
\begin{proof}
Write $Z(t)\defeq\Ztor_k(t;h)$, $I(t)\defeq I^{(r')}_k(t;h)$, $\rho\defeq\rho^{(r')}_k(t;h)$, $B\defeq B_k$, $b\defeq b_k(h)$. The bound on $b$ is \Cref{rot:prop:H6S1,rot:prop:H6S2}. By \Cref{rot:prop:H6S1} at $t$ and at $0$: $Z(t)\le e^xI(t)+B$ and $Z(0)\ge e^{-x}I(0)$ give $Z(t)/Z(0)\le e^{2x}\rho+B/Z(0)\le e^{2x}\rho+b$; and $Z(t)\ge e^{-x}I(t)=e^{-x}\rho I(0)$ with $I(0)\ge e^{-x}(Z(0)-B)$ give $Z(t)/Z(0)\ge e^{-2x}\rho(1-B/Z(0))\ge e^{-2x}(1-b)\rho$. Hence, by \Cref{rot:lem:H6A}, $\rho\le e^{2x}(1-b)^{-1}Z(t)/Z(0)\le2e^{2x}e^{C_Gt^2}\le4e^{C_Gt^2}$. For the difference: $Z(t)/Z(0)-\rho\le(e^{2x}-1)\rho+b$ and $\rho-Z(t)/Z(0)\le\rho[(1-e^{-2x})+b]\le\rho[(e^{2x}-1)+b]$; the factor $(1+4e^{C_Gt^2})$ on $b$ covers both sides. The last statement is the case $r'=r(k)$, $h\in\{h_k,\hinf\}$, with the identities of \Cref{rot:lem:H6conv}.
\end{proof}
